\section{Marco unificado para modelos generativos: de Latent a Data}

\subsection{Idea central: aprender mapas entre distribuciones}

\begin{importantbox}{Enfoque común en todos los modelos generativos}
Todos los modelos generativos profundos (GAN, VAE, Diffusion Model, Flow Model) comparten la misma idea central: \textbf{aprender un mapa desde una distribución simple (latent distribution) hacia una distribución de datos compleja (data distribution)}.

En particular:
\begin{enumerate}
    \item Seleccionar una distribución simple fácil de muestrear como latent distribution $p(z)$ (generalmente una distribución gaussiana estándar $\mathcal{N}(\mathbf{0}, \mathbf{I})$)
    \item Aprender un mapa mediante una red neuronal $G_\theta: z \mapsto x$
    \item Durante la generación: muestrear $z \sim p(z)$, calcular $x = G_\theta(z)$
\end{enumerate}
La diferencia entre los distintos modelos radica en \textbf{cómo entrenar este mapa}.
\end{importantbox}

\subsection{Convenciones de notación}

Este curso utiliza las siguientes notaciones que se mantienen a lo largo del texto:

\begin{table}[H]
\centering
\begin{tabular}{lll}
\toprule
Símbolo & Significado & Comentario \\
\midrule
$z$ & variable latent & proviene de una distribución simple (generalmente $\mathcal{N}(\mathbf{0}, \mathbf{I})$) \\
$x$ & punto de datos & proviene del conjunto de datos de entrenamiento \\
$p(z)$ & distribución latent / prior & conocida \\
$p(x)$ / $p_{\text{data}}(x)$ & distribución de datos & \textbf{desconocida}; este es el objetivo que queremos aproximar \\
$p_\theta(x|z)$ & distribución likelihood & distribución condicional para generar $x$ dado $z$ \\
$p(z|x)$ & distribución posterior & distribución condicional para inferir $z$ dado $x$ \\
\bottomrule
\end{tabular}
\caption{Convenciones de notación utilizadas en este curso}
\end{table}

\subsection{La inspiración del Autoencoder}

El autoencoder es un punto de partida para comprender los modelos generativos:

\begin{itemize}
    \item \textbf{Encoder}: $f_\phi: x \mapsto z$, comprime el punto de datos hacia un espacio latent de baja dimensión
    \item \textbf{Decoder}: $g_\theta: z \mapsto \hat{x}$, reconstruye los datos a partir de la representación latent
    \item Objetivo del entrenamiento: minimizar el error de reconstrucción $\|x - g_\theta(f_\phi(x))\|^2$
\end{itemize}

En el contexto de los modelos generativos, el \textbf{decoder es el mapa que necesitamos}; mapea puntos del espacio latent hacia el espacio de datos. Pero la pregunta clave es: \textbf{cómo garantizar que el decoder pueda mapear cualquier punto extraído aleatoriamente de la latent distribution en un punto de datos razonable}?

\begin{warningbox}{Un autoencoder no es un modelo generativo}
El espacio latent de un autoencoder estándar no sigue alguna distribución simple conocida. Si tomamos directamente un punto $z$ al azar del espacio latent y lo introducimos en el decoder, la salida resultante probablemente no parecerá datos reales. VAE resuelve este problema imponiendo restricciones previas al espacio latent.
\end{warningbox}

\subsection{Resumen de la sección}

El paradigma unificado de los modelos generativos es: partiendo de una distribución simple conocida $p(z)$, aprender mediante redes neuronales el mapa hacia la distribución de datos $p_{\text{data}}(x)$. La única diferencia entre GAN, VAE, Diffusion Model y Flow Model radica en sus estrategias de entrenamiento. El decoder del autoencoder proporciona un prototipo de este mapa, pero requiere mecanismos adicionales para garantizar la estructura del espacio latent.

\newpage
